\subsection{Induced representations}

We keep the notations and conventions of the preceding number concerning the measures \(\beta\) on H and \(\nu\) on G/H, as well as the function \(\kappa\colon\mathrm{G}\to\mathbf{R}_{+}^{*}\) .

There exists a continuous function \(\eta\) of \(\mathrm{G} \times \mathrm{G} / \mathrm{H}\) in \(\mathbf{R}_{+}^{*}\) such that

\[
\eta (x, y \mathrm{H}) = \frac {\kappa (x y)}{\kappa (x)}
\]

for all \((x,y)\in\mathrm{G}\times\mathrm{G}\) , and \(\boldsymbol{\gamma}_{\mathrm{G}/\mathrm{H}}(x)\nu=(y\mapsto\eta(x^{-1},y))\cdot\nu\) for \(x\in G\) (INT, VII, p. 56, § 2, no. 5, th. 2, c)).

Let \(p \in [1, +\infty[\) and let \(\pi\) a unitary representation of H in a complex Hilbert space E.

Lemma 6. --- Let \(f \in \mathcal{K}_{\overline{\pi}}(\mathrm{G})\) . For every \(g \in \mathrm{G}\) , the function

\[
\widetilde {f} \colon x \mapsto \eta (g ^ {- 1}, x \mathrm{H}) ^ {1 / p} f (g ^ {- 1} x)
\]

of G in E belongs to \(\mathcal{K}_{\overline{\pi}}(\mathrm{G})\) and satisfies \(\mathrm{N}_p(\widetilde{f}) = \mathrm{N}_p(f)\) .

One verifies without difficulty that \(\tilde{f} \in \mathcal{K}_{\overline{\pi}}(\mathrm{G})\) . Since

\[
\mathrm{N} _ {p} (\widetilde {f}) ^ {p} = \int_ {\mathrm{G/H}} ^ {*} \| \widetilde {f} \| ^ {p} d \nu = \int_ {\mathrm{G/H}} ^ {*} \gamma_ {\mathrm{G/H}} (g) (\| f \| ^ {p}) (y) \eta (g ^ {- 1}, y) d \nu (y)
\]

 and that \((y\mapsto \eta (g^{-1},y))\cdot \nu = \pmb{\gamma}_{\mathrm{G / H}}(g)\nu\) , we obtain \(\mathrm{N}_p(\widetilde{f}) = \mathrm{N}_p(f)\) .

 It follows from this lemma that there exists a continuous and isometric representation \(\widetilde{\pi}\) from G into \(\mathrm{L}_{\overline{\pi}}^{p}(\mathrm{G})\) such that for \(f\in \mathcal{K}_{\pi}(\mathrm{G})\) and \(g\in \mathrm{G}\) , the element \(\widetilde{\pi}(g)f\) is the class of the function \(\widetilde{f}\) defined above. If \(p=2\) , then this representation is unitary.

DEFINITION 2. --- One says that the unitary representation of G in the space \(L_{\overline{\pi}}^{2}(G)\) thus defined is the unitary representation of G induced by the representation \(\pi\) of H with respect to \(\kappa\) . We denote it by \(\operatorname{Ind}_{\mathrm{H}}^{\mathrm{G}}(\pi, \kappa)\) , or simply \(\operatorname{Ind}_{\mathrm{H}}^{\mathrm{G}}(\pi)\) .

 Remarks. --- 1) Let \(\varrho\) a unitary representation of H in a complex Hilbert space F, and let \(u\colon\pi\to\varrho\) an H-morphism. For every function \(f\in\mathcal{K}_{\overline{\pi}}(\mathrm{G})\) , denote \(v(f)\) the function \(g\mapsto u(f(g))\) from G to F; it belongs to \(\mathcal{K}_{\overline{\varrho}}(\mathrm{G})\) and satisfies \(\mathrm{N}_{p}(v(f))\leqslant\|u\|\mathrm{N}_{p}(f)\) . The linear mapping from \(\mathcal{K}_{\overline{\pi}}(\mathrm{G})\) in \(\mathcal{K}_{\overline{\varrho}}(\mathrm{G})\) which to \(f\) associates \(v(f)\) extends therefore to a continuous H-morphism from \(\mathrm{Ind}_{\mathrm{H}}^{\mathrm{G}}(\pi)\) in \(\mathrm{Ind}_{\mathrm{H}}^{\mathrm{G}}(\varrho)\) which is denoted \(\mathrm{Ind}_{\mathrm{H}}^{\mathrm{G}}(u)\) . We have \(\mathrm{Ind}_{\mathrm{H}}^{\mathrm{G}}(1_{\pi})=1_{\mathrm{Ind}_{\mathrm{H}}^{\mathrm{G}}(\pi)}\) . Let \(\sigma\) a unitary representation of H, and let \(v\colon\varrho\to\sigma\) an H-morphism; we have \(\mathrm{Ind}_{\mathrm{H}}^{\mathrm{G}}(v\circ u)=\mathrm{Ind}_{\mathrm{H}}^{\mathrm{G}}(v)\circ\mathrm{Ind}_{\mathrm{H}}^{\mathrm{G}}(u)\) .

In other words, the construction which to \(\pi\) associates \(\mathrm{Ind}_{\mathrm{H}}^{\mathrm{G}}(\pi)\) and with \(u\) associates \(\mathrm{Ind}_{\mathrm{H}}^{\mathrm{G}}(u)\) is a functor from the category of unitary representations of H to that of unitary representations of G (cf. CAT, I, § 2, in preparation).

\begin{enumerate}
\def\labelenumi{\arabic{enumi})}
\setcounter{enumi}{1}
\tightlist
\item
  Let \(\kappa^{\prime}\colon \mathrm{G}\to \mathbf{R}_{+}^{*}\) such that
\end{enumerate}

\[
\frac {\kappa^ {\prime} (x h)}{\kappa^ {\prime} (x)} = \frac {\Delta_ {\mathrm{H}} (h)}{\Delta_ {\mathrm{G}} (h)} = \frac {\kappa (x h)}{\kappa (x)}
\]

 for all \((x,h)\in\mathrm{G}\times\mathrm{H}\) . Let \(\nu^{\prime}\) the quasi-invariant measure \((\kappa^{\prime}\cdot\mu)/\beta\) on G/H. The function \(\kappa^{\prime}\kappa^{-1}\) defines, by passing to the quotient, a continuous function \(\xi\colon G/H\to R_{+}^{*}\) such that \(\nu^{\prime}=\xi\cdot\nu\) . The endomorphism \(\alpha\) of \(\mathcal{K}_{\pi}(G)\) defined by \(f\mapsto(\kappa^{\prime}\kappa^{-1})^{1/p}f\) satisfies

\[
\int_ {\mathrm{G} / \mathrm{H}} ^ {*} \| f \| ^ {p} d \nu^ {\prime} = \int_ {\mathrm{G} / \mathrm{H}} ^ {*} \| \alpha (f) \| ^ {p} d \nu
\]

and makes it possible to identify the spaces \(\mathcal{L}_{\pi}^{p}(\mathrm{G},\nu^{\prime})\) and \(\mathcal{L}_{\pi}^{p}(\mathrm{G},\nu)\) , as well as the spaces \(\mathrm{L}_{\pi}^{p}(\mathrm{G},\nu^{\prime})\) and \(\mathrm{L}_{\pi}^{p}(\mathrm{G},\nu)\) . Moreover, \(\alpha\) induces an isometric isomorphism of representations \(\mathrm{Ind}_{\mathrm{H}}^{\mathrm{G}}(\pi,\kappa^{\prime})\) and \(\mathrm{Ind}_{\mathrm{H}}^{\mathrm{G}}(\pi,\kappa)\) .

